\noindent\textbf{(a)} $\mu_1=(y^2-x^2+1)^{-2}$ y $x/(y^2-x^2+1)=C$ donde el denominador no se anula. Se recupera separadamente $y^2=x^2-1$.

\noindent\textbf{(b)} $\mu_2=x^{-2}$ y $y^2=x^2+Kx-1$. Las soluciones gráficas reales son $y=\pm\sqrt{x^2+Kx-1}$ en intervalos de radicando estrictamente positivo.

\noindent\textbf{(c)} La familia completa es $y=\pm\sqrt{x^2+Kx-1}$, con radicando positivo. El caso $K=0$ es precisamente el que excluye el primer factor.
